\subsection{Why the coverage is what it is}

The coverage of the anchored interval is high by construction. With a second source that counts what the transition page would count if it followed every \wdn{} rating, Proposition~\ref{prop:containment} shows that the interval contains the Jeffreys interval of the oracle, so its coverage cannot be below the oracle's and its width is close to twice the oracle's, the ratio that the numerical check on random cohorts and the experiments both give. In the first-year and ten-fold designs the coverage experiment is thus not a test that the method could have failed; in the pooled design the oracle itself covers \nv{cov.pool.i9.followed}\% and the experiment could have failed, and the interval covered \nv{cov.pool.i9.anchored}\%. The Manski band covers as often or more often in every design, so coverage does not separate the anchored interval from the band; its width does, by a factor of \nv{wr.fy.i9} to \nv{wr.x10.i9}. It measures how far above the nominal level the interval sits, from \nv{cov.fy.i9.anchored}\% at informativeness 9 in the first annual cohort to \nv{cov.x10.i9.anchored}\% in the cohort ten times larger, and what it costs in width. A validator who wants an interval at its nominal level should not use this one; one who wants an interval that does not mislead about what the transition page can say should.

The experiments that could fail the method are those that stress the second source, and they show where it is weak. The coverage survives a second source that counts half or 1.3 times the defaults, as the interval is wide. The centre does not: its bias is that of the capture, and the window in which anchoring beats removal on bias is narrow, from $\rho$ to $2-\rho$ for a removal ratio $\rho$ of \nv{rho.fy.i9}. A validator who knows that the second source counts more events than the transition page, as the default page does, knows the sign of the bias of the anchored point.

\subsection{What the agency cohorts say}

The pooled anchored $k$ of S\&P's CCC to C is \nv{real.sp.CCC-C.k}, and that of grade B \nv{real.sp.B.k}. Taken at face value, withdrawal in these grades would be an order of magnitude more informative than independence, and the PD of grade B would be \nv{real.sp.B.pdanc}\% rather than \nv{real.sp.B.pdind}\%. The stability experiment says that this reading is not warranted, which is consistent with the finding of \citet{cantor2007} that the data they examine support random censoring: an anchored $k$ far above 1 that does not hold across years is better read as the difference between two definitions of default than as informative withdrawal. On known truth, an exact second source gives annual bands that hold the pooled $k$ in \nv{stab.i9.c1.B.share}\% of the cohorts, and a definition gap drawn each year over the range observed for grade B, \nv{ratio.B.min} to \nv{ratio.B.max} times the observed defaults and centred above the observed ratios, brings the share to \nv{stab.i9.cdobs.B.share}\%, near the \nv{real.sp.B.share}\% observed for S\&P's grade B (1.0 to 4.0 times the followed defaults gives \nv{stab.i9.cu1.0-4.0.B.share}\%); drawn over the range observed for CCC to C it leaves the share at \nv{stab.i9.cdobs.CCC-C.share}\%, against \nv{real.sp.CCC-C.share}\% observed. The anchored $k$ of the agencies' cohorts absorbs the difference between the two pages' definitions of default, which changes from year to year, together with whatever informativeness the withdrawals have. It is an anchored $k$ conditional on the two definitions being equivalent, and the repository does not make them so. The case of Fitch is the sharpest: its transition pages hold no defaulted rating in \nv{fitch.empty.n} years while its default page counts \nv{fitch.empty.lost}, so the anchored $k$ of Fitch is mostly a statement about where the defaulted ratings are recorded.

What the agency cohorts do support is the statement that the tipping $k$ matters. For S\&P's grade B, the tipping $k$ to the next grade is \nv{real.sp.B.tip}, and the removal PD would need withdrawal to be that informative to move the grade to the next notch; the anchored $k$, whatever its meaning, is below it. For Fitch's grade BB, the anchored $k$ is \nv{real.fitch.BB.k} against a tipping $k$ of \nv{real.fitch.BB.tip}, and the decision flips under the anchor but only if the definitions agree.

\subsection{Negative results and what they imply}

The prediction as a whole is refuted, by its stability clause on the agencies' data and by its coverage clause against the Manski band. H2 is refuted in part: the clause that removal and keeping cover under 50\% of the repetitions at informativeness 9 fails in the first annual cohort, \nv{cov.fy.i9.industry}\% and \nv{cov.fy.i9.unfollowed}\%, because the bias of about a fifth to a quarter is smaller than the sampling interval of a grade with some sixteen defaults a year. It holds when the cohort is large enough or the years are pooled. H3 fails as worded, because the Manski band covers as often or more often than the anchored interval, which buys width and not coverage. H4 is refuted: the clause on stability fails on the agencies' cohorts. The continuous-time estimator that reads withdrawal as independent censoring has a bias between those of the pooled removal and keeping estimates at informativeness 9, $-\nv{dur.i9.abias}\%$ against $\nv{bias.pool.i9.industry}\%$ and $\nv{bias.pool.i9.unfollowed}\%$, which shows that the problem is not removed by working in continuous time.

For a bank the implications are three. A grade's PD under withdrawal should be reported as the family and its tipping $k$ at the thresholds that matter, which needs no second source and no assumption beyond A1 and A3. The anchored interval should be reported only where the definitions of the two sources have been compared, and with the capture window of Proposition~\ref{prop:capture} stated. And pooling years, the usual practice, brings a selection bias that no estimator of the cohort corrects, and that all the intervals of the paper inherit.
